\documentclass[11pt]{article}

\usepackage[margin=1in]{geometry}
\usepackage{amsmath}
\usepackage{amssymb}
\usepackage{enumitem}
\usepackage{hyperref}
\usepackage{booktabs}

\title{Architecture Note\\Efficient Implementation and Comparative Analysis of Graph Centrality Measures}
\author{Project 19}
\date{\today}

\begin{document}
\maketitle

\section{System Objective}
The system provides an end-to-end workflow for graph centrality analysis on unweighted undirected graphs, including algorithm implementation, correctness validation, benchmarking, ranking comparison, plotting, and interactive visualization.

\section{Module Overview}
Core modules under \texttt{src/} are organized by responsibility:
\begin{itemize}[leftmargin=*]
    \item \texttt{graph\_utils.py}: adjacency-list graph model, edge-list loader, synthetic generators, graph summaries, connectivity checks.
    \item \texttt{closeness.py}: BFS shortest paths and exact closeness centrality.
    \item \texttt{betweenness.py}: Brandes algorithm for exact betweenness centrality.
    \item \texttt{eigenvalue.py}: power iteration for eigenvalue centrality with convergence metadata.
    \item \texttt{validators.py}: structural and ranking sanity checks, optional tiny-graph comparisons against NetworkX.
    \item \texttt{experiments.py}: synthetic benchmark suite with runtime collection and export.
    \item \texttt{compare\_rankings.py}: top-$k$ overlap, Jaccard, Spearman, and exportable comparison summaries.
    \item \texttt{visualize.py}: generation of runtime, distribution, overlap, correlation, and graph highlight plots.
    \item \texttt{demo.py}: integrated command-line demo for centrality outputs.
\end{itemize}

Supporting entry points:
\begin{itemize}[leftmargin=*]
    \item \texttt{run\_project.py}: centrality demonstration run.
    \item \texttt{run\_experiments.py}: benchmark execution and export.
    \item \texttt{run\_analysis.py}: ranking comparison and plot pipeline.
\end{itemize}

Dashboard modules under \texttt{dashboard/}:
\begin{itemize}[leftmargin=*]
    \item \texttt{app.py}: Streamlit application layout and interactions.
    \item \texttt{viz\_utils.py}: dashboard graph parsing, metric computation, table export, and in-app rendering support.
\end{itemize}

\section{Data Flow}
The system follows a staged flow:
\begin{enumerate}[leftmargin=*]
    \item \textbf{Graph acquisition}: synthetic generation or edge-list parsing.
    \item \textbf{Metric computation}: closeness, betweenness, eigenvalue centrality.
    \item \textbf{Validation}: deterministic sanity checks and optional micro-comparisons.
    \item \textbf{Benchmarking}: repeated metric execution across graph families and sizes.
    \item \textbf{Comparison analysis}: overlap and correlation summaries between ranking outputs.
    \item \textbf{Visualization}: static plots for runtime and ranking behavior.
    \item \textbf{Interactive exploration}: dashboard-based ad hoc analysis and CSV export.
\end{enumerate}

\section{Graph Representation and API}
Graphs are stored using adjacency lists with contiguous internal node IDs and bidirectional mappings between original labels and internal IDs. This design balances algorithmic efficiency and user-facing interpretability.

\section{Algorithm Integration}
\subsection{Closeness Centrality}
For node $v$, with reachable set $R(v)$ and $r = |R(v)|$:
\[
S(v)=\sum_{u \in R(v)} d(v,u), \qquad C_{\text{base}}(v)=\frac{r}{S(v)}.
\]
The implementation uses BFS from each source and applies reachable-fraction normalization:
\[
C(v)=C_{\text{base}}(v)\cdot\frac{r}{n-1}.
\]

\subsection{Betweenness Centrality}
The implementation follows Brandes for unweighted undirected graphs. For each source $s$, BFS builds shortest-path predecessor sets and path counts $\sigma$. Reverse traversal accumulates dependencies and aggregates node contributions. Undirected double counting is corrected, and optional normalization is applied:
\[
\frac{2}{(n-1)(n-2)}.
\]

\subsection{Eigenvalue Centrality}
Power iteration estimates the principal eigenvector:
\[
x^{(k+1)}=\frac{(A+\alpha I)x^{(k)}}{\|(A+\alpha I)x^{(k)}\|_2},
\]
where a small shift $\alpha$ improves convergence stability for bipartite structures.

\section{Benchmarking Pipeline}
The benchmark runner executes each metric across synthetic graph families with controlled sizes and fixed seeds where randomization is used. For each graph/metric pair, it stores:
\begin{itemize}[leftmargin=*]
    \item graph family and configuration,
    \item node and edge counts,
    \item connectivity,
    \item runtime measured by \texttt{time.perf\_counter},
    \item status and metric-specific notes.
\end{itemize}
Exports are written to CSV, JSON, and timestamped log files under \texttt{outputs/}.

\section{Analysis Pipeline}
The analysis runner consumes benchmark output and computed centrality vectors to produce:
\begin{itemize}[leftmargin=*]
    \item pairwise top-$k$ overlap and Jaccard similarity,
    \item pairwise Spearman rank correlations,
    \item summary exports to \texttt{outputs/comparison\_results.csv} and \texttt{outputs/comparison\_results.json},
    \item plot assets in \texttt{outputs/plots/} and copied report figures in \texttt{report/figures/}.
\end{itemize}

\section{Dashboard Structure}
The Streamlit dashboard supports:
\begin{itemize}[leftmargin=*]
    \item synthetic graph selection and configurable parameters,
    \item edge-list upload and parsing,
    \item graph summary metrics,
    \item selected metric computation,
    \item node-level visualization with centrality-based encoding,
    \item top-$k$ ranking and pairwise comparison display,
    \item CSV export of computed metric scores.
\end{itemize}

\section{Limitations}
\begin{itemize}[leftmargin=*]
    \item Primary algorithmic support is restricted to unweighted undirected graphs.
    \item Benchmark workloads are synthetic and moderate in size by design.
    \item Dashboard interaction is tuned for small-to-medium graph sizes, not large-scale interactive rendering.
    \item Advanced statistical significance and uncertainty analysis across rankings are not yet integrated.
\end{itemize}

\section{Planned Extensions}
\begin{itemize}[leftmargin=*]
    \item weighted and directed graph centrality variants,
    \item larger real-world dataset ingestion and curation,
    \item richer comparison statistics and stability analyses,
    \item enhanced dashboard support for larger graph exploration.
\end{itemize}

\end{document}
